\documentclass[letter]{core/polinetwork}

\pnTitle{Cos'è PoliNetwork}

\begin{document}

\pnMakeHeader

% === CHI SIAMO ===
\section{Chi siamo}

PoliNetwork è la più grande organizzazione studentesca indipendente
legata al Politecnico di Milano. Siamo un'Associazione di Promozione
Sociale (APS) no-profit, iscritta come Ente del Terzo Settore,
gestita interamente da studenti e per studenti.

La nostra missione è semplice: \textbf{rendere più semplice la vita
	universitaria al Politecnico di Milano}. Lo facciamo attraverso
strumenti digitali, community, eventi e supporto diretto, colmando
le lacune comunicative e burocratiche che ogni studente incontra nel
proprio percorso.

PoliNetwork è \textbf{indipendente dal Politecnico di Milano}: non
siamo un'associazione dell'Albo istituzionale dell'Ateneo.
Questa indipendenza è una scelta deliberata che ci permette di agire
con flessibilità, sviluppare infrastrutture autonome e rispondere
alle esigenze reali degli studenti senza vincoli istituzionali.

Tutti i nostri progetti di carattere tecnico-informatico sono
rilasciati con \textbf{licenza open source}, come previsto dal
nostro Statuto.

% === COSA FACCIAMO ===
\section{Cosa facciamo}

\subsection{La rete di gruppi}

Gestiamo \textbf{centinaia di gruppi WhatsApp e Telegram}, organizzati
per corso di laurea, scaglione, sede e argomento. Ogni anno, decine di
migliaia di studenti — dalle matricole ai laureandi — utilizzano i
nostri gruppi per:

\begin{enumerate}[leftmargin=2em]
	\item Scambiarsi appunti e informazioni sugli esami
	\item Ricevere risposte rapide su burocrazia, ISEE, borse di studio, piano degli studi
	\item Trovare alloggi, libri usati, ripetizioni
	\item Connettersi con altri studenti del proprio corso
\end{enumerate}

I gruppi sono moderati da una rete di \textbf{volontari (admin)} che
garantiscono qualità, rispetto delle regole e assenza di spam.

\subsection{Strumenti digitali}

Sviluppiamo e manteniamo strumenti gratuiti utilizzati da migliaia
di studenti:

\begin{enumerate}[leftmargin=2em]
	\item \textbf{Graduatorie TOL}: il sistema che rende consultabili
	      e navigabili le graduatorie dei Test Online del Politecnico
	\item \textbf{Simulatore TOL}: per prepararsi al test d'ingresso
	\item \textbf{Bot Telegram}: per trovare rapidamente i gruppi del
	      proprio corso
	\item \textbf{Repository di materiale didattico}: archivi condivisi
	      con dispense, appunti e temi d'esame
	\item \textbf{Sito web}: hub centrale con guide per le matricole,
	      mappe, FAQ e risorse
\end{enumerate}

\subsection{Eventi}

Organizziamo eventi di aggregazione per la community studentesca:
aperitivi, welcome party per le matricole, eventi interni per i
nostri membri e collaborazioni con partner esterni. Partecipiamo
anche a festival universitari e Open Day con stand dedicati.

\subsection{Partnership}

Collaboriamo con altre realtà dell'ecosistema universitario per
creare valore per gli studenti: da eventi congiunti a iniziative
di visibilità reciproca, sempre nel rispetto della nostra
indipendenza e dei nostri valori.

% === I NOSTRI TEAM ===
\section{I nostri team}

PoliNetwork è organizzata in \textbf{quattro team operativi},
ciascuno con un proprio ambito di responsabilità. Entrare in un
team significa contribuire attivamente alla crescita
dell'associazione e acquisire competenze concrete.

\subsection{IT}
Gestisce tutta l'infrastruttura tecnica di PoliNetwork: dal
sito web agli strumenti per gli studenti (graduatorie, simulatore,
bot). Sviluppa nuovi progetti open source. Richiede competenze
di programmazione.

\subsection{HR}
Si occupa del reclutamento e della gestione delle persone:
colloqui per nuovi admin e membri dei team, censimento degli
admin attivi, organizzazione di eventi interni in collaborazione
con il team eventi.

\subsection{Design \& Marketing}
Cura la comunicazione dell'associazione su tutti i canali: gestione
dei social media, creazione di contenuti grafici, strategia di
branding e visibilità. Cerchiamo sia persone con competenze
grafiche (Figma, Canva, ecc.) sia persone abili nella scrittura
e nella strategia comunicativa.

\subsection{Events \& Partnerships}
Organizza e gestisce gli eventi di PoliNetwork (esterni, interni,
in collaborazione). Si occupa anche di avviare e mantenere
partnership con aziende, altre associazioni e realtà esterne
all'ateneo.

\vspace{1em}

Oltre ai team, PoliNetwork si regge su una rete di
\textbf{admin volontari} che moderano i gruppi quotidianamente.
Diventare admin è un altro modo per contribuire e far parte della
community.

% === FAR PARTE DI POLINETWORK ===
\section{Far parte di PoliNetwork}

\subsection{Perché unirti a noi}

\begin{enumerate}[leftmargin=2em]
	\item \textbf{Crescita personale}: lavorerai su progetti reali
	      fin dal primo giorno, non simulazioni
	\item \textbf{Competenze trasferibili}: organizzazione eventi,
	      gestione team, comunicazione, sviluppo software, recruiting
	      — tutto spendibile nel mondo del lavoro
	\item \textbf{Network}: conoscerai centinaia di studenti
	      motivati da tutti i corsi del Politecnico
	\item \textbf{Impatto concreto}: i servizi che costruirai
	      verranno utilizzati da decine di migliaia di persone
	\item \textbf{Flessibilità}: non abbiamo una sede fisica e gran
	      parte del lavoro è da remoto, con incontri periodici
	      dal vivo
\end{enumerate}

\subsection{Cosa cerchiamo}

Cerchiamo persone \textbf{proattive}, che non aspettino che qualcuno
dica loro cosa fare, ma che propongano, sperimentino e si prendano
responsabilità. Non servono competenze pregresse per tutti i ruoli
— serve voglia di fare e di imparare.

\subsection{Come candidarsi}

Il processo di selezione prevede:
\begin{enumerate}[leftmargin=2em]
	\item \textbf{Candidatura}: compilazione del form disponibile
	      sui nostri canali
	\item \textbf{Colloquio conoscitivo}: una conversazione informale
	      per conoscerti e capire le tue motivazioni
	\item \textbf{Eventuale colloquio specifico}: per i ruoli nei
	      team, un secondo colloquio con una persona del team
	      di riferimento
	\item \textbf{Periodo di inserimento}: i primi mesi servono a
	      entrambi per capire se il match funziona
\end{enumerate}

Le candidature per i team sono aperte durante i periodi di
recruiting comunicati sui nostri canali social. Se desideri
diventare admin, puoi contattarci in qualsiasi momento.

\subsection{Diventare socio}

I membri più attivi possono diventare \textbf{soci} dell'associazione.
La quota associativa è simbolica (5\euro{}/anno) e include:
\begin{enumerate}[leftmargin=2em]
	\item Email istituzionale \texttt{nome.cognome@polinetwork.org}
	\item 1 TB di spazio OneDrive
	\item Diritto di voto nelle Assemblee dei Soci
	\item Gadget e lanyard dell'associazione
\end{enumerate}

% === CONTATTI ===
\section{Contatti e link utili}

\begin{enumerate}[leftmargin=2em]
	\item \textbf{Sito web}: \pnWebsiteLink
	\item \textbf{Instagram}: \pnInstagramLink
	\item \textbf{Email}: \href{mailto:info@polinetwork.org}{\textcolor{pn-primary}{\texttt{info@polinetwork.org}}}
	\item \textbf{Telegram}: cerca \pnTelegramLink\ per
	      trovare i gruppi del tuo corso
\end{enumerate}

\end{document}
